% Main template for theory notes and problem solutions.
\documentclass[a4paper, 12pt]{article}
% Shared preamble: packages, boxes, commands. NOT a standalone document.

\usepackage{amsthm, amsfonts, amssymb, mathtools}
\usepackage[utf8]{inputenc}
\usepackage[T2A]{fontenc}
\usepackage[english, russian]{babel}
\usepackage{microtype}
\usepackage{tikz}
\usepackage[most]{tcolorbox}
\usepackage[left=2cm,right=2cm,top=2cm,bottom=2cm]{geometry}
\usepackage{graphicx}
\usepackage{caption}
\usepackage{xcolor}
\usepackage{array}
\usetikzlibrary{shapes.symbols}
% РЕШЕНИЕ ПРОБЛЕМЫ С color И babel
\usepackage{etoolbox}
\makeatletter
\patchcmd{\endtcolorbox}{\end{lrbox}}{\color@reset\end{lrbox}}{}{}
\patchcmd{\tcb@break@end}{\end{tcb@savebox}}{\color@reset\end{tcb@savebox}}{}{}
\makeatother

\widowpenalty=10000

\newtcolorbox{dashframe}{
  enhanced,
  frame hidden, % полностью убираем стандартную рамку
  interior hidden,
  arc=2mm,
  left=6pt, right=6pt, top=6pt, bottom=6pt,
  breakable,
  overlay={
    \draw[dash pattern=on 4pt off 2pt]
      (frame.south west) rectangle (frame.north east);
  }
}

\newtcolorbox{corollary}[1][]{
  enhanced,
  colback=cyan!5,
  colframe=cyan!60,
  breakable,
  fonttitle=\bfseries,
  left=8pt,
  right=8pt,
  top=10pt,
  bottom=8pt,
  title={%
    \textbf{Следствие}%
    \IfBlankTF{#1}{}{ \text{#1}}%
  }
}

\newcounter{theorem}[section]

\newtcolorbox[auto counter, number within=section]{theorem}[1][]{
  enhanced,
  colback=blue!4,
  colframe=blue!60!black,
  fonttitle=\bfseries,
  left=8pt,
  right=8pt,
  top=10pt,
  bottom=8pt,
  breakable,
  title={%
    \textbf{Теорема \thetcbcounter}%
    \IfBlankTF{#1}{}{ ({\itshape #1})}%
  }
}
\newtcolorbox{thmproof}{
  enhanced,
  breakable,
  colback=gray!3,
  colframe=gray!50!black,
  coltitle=black,
  fonttitle=\bfseries\itshape,
  left=17pt,
  right=8pt,
  top=8pt,
  bottom=8pt,
  boxrule=0.8pt,
  arc=4pt,
}

\newcommand{\circled}[1]{%
  \tikz[baseline=(char.base)]{%
    \node[circle, draw, minimum size=0.3em, inner sep=2pt, line width=0.06em] (char) {%
      \normalfont #1%
    };
  }%
}

\newcommand{\cloud}[1]{%
  \begin{tikzpicture}
    \node[cloud, draw, fill=white, aspect=3, inner sep=3pt] {#1};
  \end{tikzpicture}%
}

\linespread{1.3}

\tcbuselibrary{theorems, skins}

\newtcolorbox[auto counter, number within=section]{lemma}[1][]{
  enhanced,
  breakable,
  colback=green!5,
  colframe=green!60!black,
  fonttitle=\bfseries,
  left=8pt,
  right=8pt,
  top=10pt,
  bottom=8pt,
  title={%
    \textbf{Лемма \thetcbcounter}%
    \IfBlankTF{#1}{}{ ({\itshape #1})}%
  }
}

\newtcolorbox{definition}[1][]{
  enhanced,
  breakable,
  colback=orange!5,
  colframe=orange!60!black,
  fonttitle=\bfseries,
  left=8pt,
  right=8pt,
  top=10pt,
  bottom=8pt,
  title={%
    \textbf{Определение}%
    \IfBlankTF{#1}{}{ ({\itshape #1})}%
  }
}

\newtcolorbox{designation}[1][]{
  breakable,
  enhanced,
  colback=purple!5,
  colframe=purple!60!black,
  fonttitle=\bfseries,
  left=8pt,
  right=8pt,
  top=10pt,
  bottom=8pt,
  title={%
    \textbf{Обозначение}%
    \IfBlankTF{#1}{}{ ({\itshape #1})}%
  }
}

\newtcolorbox{example}[1][]{
  breakable,
  enhanced,
  colback=teal!5,
  colframe=teal!60!black,
  fonttitle=\bfseries,
  left=8pt,
  right=8pt,
  top=10pt,
  bottom=8pt,
  title={%
    \textbf{Пример}%
    \IfBlankTF{#1}{}{ ({\itshape #1})}%
  },
}

\newtcolorbox{nota bene}[1][]{
  breakable,
  enhanced,
  colback=red!5,
  colframe=red!60!black,
  fonttitle=\bfseries,
  left=8pt,
  right=8pt,
  top=10pt,
  bottom=8pt,
  title={%
    \textbf{\circled{NB}}%
    \IfBlankTF{#1}{}{ ({\itshape #1})}%
  }
}

\newcolumntype{C}[1]{>{\centering\let\newline\\\arraybackslash\hspace{0pt}}m{#1}}
%TeXporter:DefineEmojiColors
\definecolor{c0a0808}{RGB}{10,8,8}
\definecolor{cff0000}{RGB}{255,0,0}
\definecolor{cff9700}{RGB}{255,151,0}
\definecolor{c1a0f00}{RGB}{26,15,0}
\definecolor{c782121}{RGB}{120,33,33}
\definecolor{cffffff}{RGB}{255,255,255}
\definecolor{cffefef}{RGB}{255,239,239}
\definecolor{c250303}{RGB}{37,3,3}
\definecolor{ce3dbdb}{RGB}{227,219,219}
\definecolor{cd45500}{RGB}{212,85,0}
\definecolor{cffd42a}{RGB}{255,212,42}
\definecolor{caa4400}{RGB}{170,68,0}
\definecolor{cff5555}{RGB}{255,85,85}
\definecolor{cffff00}{RGB}{255,255,0}
\definecolor{c5a5a5a}{RGB}{90,90,90}
\definecolor{ce9c900}{RGB}{233,201,0}
\definecolor{c2b0000}{RGB}{43,0,0}
\definecolor{c00ffff}{RGB}{0,255,255}
%TeXporter:EndDefineEmojiColors

\newcommand{\CreateSadFace}{%
\begin{tikzpicture}[scale=0.25]  % изменяйте масштаб (0.25 = маленький)
% Head
\pgfdeclareradialshading{ballshading}{\pgfpoint{-10bp}{10bp}}
{color(0bp)=(yellow!40!white);
color(9bp)=(orange!85!white);
color(18bp)=(orange!90!blue);
color(25bp)=(yellow!80!blue);
color(50bp)=(black)}
\pgfpathcircle{\pgfpoint{0cm}{0cm}}{2cm}
\pgfshadepath{ballshading}{20}
\pgfusepath{}
\pgfdeclareradialshading{ballshading}{\pgfpoint{-10bp}{10bp}}
{color(0bp)=(yellow!40!white);
color(9bp)=(yellow!30!white);
color(18bp)=(yellow!60!orange);
color(25bp)=(yellow!50!orange);
color(50bp)=(orange)}
\pgfpathcircle{\pgfpoint{0cm}{0cm}}{1.85cm}
\pgfshadepath{ballshading}{20}
\pgfusepath{}
% Head hat
\path[yshift=138,xshift=-140,scale=0.012,yscale=-1, fill=white,opacity=0.234] (411.4196,246.6473) .. controls
(344.7661,246.6473) and (288.7870,289.5501) .. (273.0134,347.5848) .. controls
(311.3176,336.0898) and (355.6292,329.4911) .. (402.8571,329.4911) .. controls
(457.9662,329.4911) and (509.1043,338.4615) .. (551.3571,353.7411) .. controls
(537.9844,292.6125) and (480.4076,246.6473) .. (411.4196,246.6473) -- cycle;

\draw[inner color=white ,outer color=c782121!20!white]
(-0.8,0.1) circle (0.35);
\fill[inner color=black ,outer color=black!20!white]
(-0.8,-0.05) circle (0.15);

\draw[inner color=white ,outer color=c782121!20!white]
(0.8,0.1) circle (0.35);
\fill[inner color=black ,outer color=black!20!white]
(0.8,-0.05) circle (0.14);

\path[xshift=-90,yshift=122,scale=0.06,yscale=-1,draw=brown!70!black,inner color=c782121!70!yellow,outer color=c782121!60!black]
(46.74375,50.94375) .. controls (44.20450,57.00941) and (38.54312,61.64973) ..
(31.56875,63.20938) .. controls (32.36042,63.35032) and (33.17901,63.42500) ..
(34.01562,63.42500) .. controls (41.06052,63.42500) and (46.77188,58.18907) ..
(46.77188,51.73125) .. controls (46.77188,51.46641) and (46.76266,51.20420) ..
(46.74375,50.94375) -- cycle;

\path[xshift=90,yshift=122,scale=0.06,yscale=-1,xscale=-1,draw=brown!70!black,inner color=c782121!70!yellow,outer color=c782121!60!black]
(46.74375,50.94375) .. controls (44.20450,57.00941) and (38.54312,61.64973) ..
(31.56875,63.20938) .. controls (32.36042,63.35032) and (33.17901,63.42500) ..
(34.01562,63.42500) .. controls (41.06052,63.42500) and (46.77188,58.18907) ..
(46.77188,51.73125) .. controls (46.77188,51.46641) and (46.76266,51.20420) ..
(46.74375,50.94375) -- cycle;

\path[xshift=55,yshift=-10,scale=0.0065,yscale=-0.8,xscale=-0.8,draw=black,inner color=cff5555!80!black,outer color=c782121!60!black]
(359.6464,67.3100) .. controls (278.1326,67.3100)
and (211.4129,129.5159) .. (205.8964,208.2475) .. controls (250.9411,190.8018)
and (304.2674,180.7163) .. (361.4277,180.7163) .. controls (417.0568,180.7163)
and (469.0358,190.2795) .. (513.3027,206.8725) .. controls (507.1094,128.7941)
and (440.6844,67.3100) .. (359.6464,67.3100) -- cycle;
\fill[xshift=55,yshift=-10,scale=0.0065,yscale=-0.8,xscale=-0.8,inner color=white,outer color=ce3dbdb!80!black,line width=0.2]
(359.6346,72.6265) .. controls (315.4332,72.6265)
and (275.6110,90.9215) .. (247.5096,120.2202) .. controls (283.3018,112.3912)
and (322.3429,108.0640) .. (363.1971,108.0640) .. controls (400.7202,108.0640)
and (436.6908,111.7186) .. (470.0408,118.3765) .. controls (442.0468,90.1452)
and (402.9104,72.6265) .. (359.6346,72.6265) -- cycle;

\fill[xshift=-63,yshift=5,scale=0.025,yscale=-1,inner color=blue!10!white,outer color=blue!80!white]
(37.38189,16.83071) -- (24.98031,58.99606) .. controls (24.04934,63.33741) and
(24.68778,67.99808) .. (26.75197,71.92913) .. controls (27.87220,74.06251) and
(29.41478,75.99641) .. (31.36401,77.41300) .. controls (33.31324,78.82959) and
(35.68099,79.70975) .. (38.09055,79.72441) .. controls (40.70723,79.74033) and
(43.28208,78.73161) .. (45.33407,77.10784) .. controls (47.38607,75.48408) and
(48.93151,73.27201) .. (49.96063,70.86614) .. controls (51.54908,67.15267) and
(51.92699,62.93268) .. (51.02362,58.99606) -- (37.38189,16.83071);

\fill[xshift=-55,yshift=-9,scale=0.017,yscale=-1,inner color=blue!5!white,outer color=blue!20!white,opacity=0.1]
(37.38189,16.83071) -- (24.98031,58.99606) .. controls (24.04934,63.33741) and
(24.68778,67.99808) .. (26.75197,71.92913) .. controls (27.87220,74.06251) and
(29.41478,75.99641) .. (31.36401,77.41300) .. controls (33.31324,78.82959) and
(35.68099,79.70975) .. (38.09055,79.72441) .. controls (40.70723,79.74033) and
(43.28208,78.73161) .. (45.33407,77.10784) .. controls (47.38607,75.48408) and
(48.93151,73.27201) .. (49.96063,70.86614) .. controls (51.54908,67.15267) and
(51.92699,62.93268) .. (51.02362,58.99606) -- (37.38189,16.83071);
\fill[xshift=-52,yshift=-17,scale=0.013,yscale=-1,fill=white,opacity=0.4]
(37.38189,16.83071) -- (24.98031,58.99606) .. controls (24.04934,63.33741) and
(24.68778,67.99808) .. (26.75197,71.92913) .. controls (27.87220,74.06251) and
(29.41478,75.99641) .. (31.36401,77.41300) .. controls (33.31324,78.82959) and
(35.68099,79.70975) .. (38.09055,79.72441) .. controls (40.70723,79.74033) and
(43.28208,78.73161) .. (45.33407,77.10784) .. controls (47.38607,75.48408) and
(48.93151,73.27201) .. (49.96063,70.86614) .. controls (51.54908,67.15267) and
(51.92699,62.93268) .. (51.02362,58.99606) -- (37.38189,16.83071);
\end{tikzpicture}%
}


% Document metadata
\title{Т5. Броуновский мост и цель bridge matching}
\author{Roman Dolgoshei}
\date{\today}

\begin{document}

\maketitle
\tableofcontents
\newpage

\section{Первое задание (Броуновский мост)}

\begin{definition}[Броуновский мост]
  Броуновским мостом из $x_0$ в $x_1$ на отрезке $[0, 1]$ называется винеровский процесс, обусловленный на $W_0 = x_0$ и $W_1 = x_1$.
\end{definition}

\textbf{Задание} \textit{(Броуновский мост)}. Вывести распределение промежуточной точки, например как условное гауссово распределение $W_t$ при известном $W_1$.

\textbf{Решение.} Рассмотрим более общий случай: мост из $x_a$ в $x_b$ на отрезке $[a,b]$. Найдём распределение $W_t$ при $a < t < b$.

\begin{equation}\label{eq:bayes}
  p(x_t \mid x_a, x_b) = \frac{p(x_b \mid x_t) p(x_t \mid x_a)}{p(x_b \mid x_a)} \qquad \text{(формула Байеса и марковость процесса)}
\end{equation}

Для винеровского процесса с дисперсией $\varepsilon$ за единицу времени, т.е. с приращениями $W_t - W_s \sim \mathcal N(0, \varepsilon (t - s) I)$:
\[
\begin{aligned}
W_t\mid W_a=x_a
&\sim \mathcal N(x_a,\varepsilon(t-a)I),\\
W_b\mid W_t=x_t
&\sim \mathcal N(x_t,\varepsilon(b-t)I),\\
W_b\mid W_a=x_a
&\sim \mathcal N(x_a,\varepsilon(b-a)I).
\end{aligned}
\]

Подставим эти плотности в формулу \eqref{eq:bayes}; через $C$ обозначим множитель, не зависящий от $x_t$:

\[
\begin{aligned}
p(x_t \mid x_a, x_b)
&= \frac{
\frac{1}{[2\pi\varepsilon(b-t)]^{d/2}}
\exp\left(-\frac{\|x_b-x_t\|^2}{2\varepsilon(b-t)}\right)
\cdot
\frac{1}{[2\pi\varepsilon(t-a)]^{d/2}}
\exp\left(-\frac{\|x_t-x_a\|^2}{2\varepsilon(t-a)}\right)
}{
\frac{1}{[2\pi\varepsilon(b-a)]^{d/2}}
\exp\left(-\frac{\|x_b-x_a\|^2}{2\varepsilon(b-a)}\right)
} \\
&= C \cdot
\frac{
\exp\left(-\frac{\|x_b-x_t\|^2}{2\varepsilon(b-t)}\right)
\exp\left(-\frac{\|x_t-x_a\|^2}{2\varepsilon(t-a)}\right)
}{
\exp\left(-\frac{\|x_b-x_a\|^2}{2\varepsilon(b-a)}\right)
} \\
&= C \,\exp\left(-\frac{\|x_b-x_t\|^2}{2\varepsilon(b-t)}-\frac{\|x_t-x_a\|^2}{2\varepsilon(t-a)}+\frac{\|x_b-x_a\|^2}{2\varepsilon(b-a)}\right) \\ &\propto \exp\left( -\frac{1}{2\varepsilon} \left[ \frac{\|x_t - x_a\|^2}{t - a} + \frac{\|x_b - x_t\|^2}{b - t} \right]\right)
\end{aligned}
\] 

Выделим полный квадрат по $x_t$:
\[
\frac{\|x_t - x_a\|^2}{t - a} + \frac{\|x_b - x_t\|^2}{b - t}
=
\frac{b-a}{(t-a)(b-t)}
\left\|x_t-\frac{(b-t)x_a+(t-a)x_b}{b-a}\right\|^2
+ \frac{\|x_b-x_a\|^2}{b-a}.
\]
Последнее слагаемое не зависит от $x_t$, поэтому плотность пропорциональна
\(
\exp\left(
-\frac{
\left\|x_t-
\frac{(b-t)x_a+(t-a)x_b}{b-a}
\right\|^2
}{
2\varepsilon\frac{(t-a)(b-t)}{b-a}
}
\right)
\), которое имеет вид $\exp\left(-\frac{\|x - \mu\|^2}{2\sigma^2}\right)$, а значит 

\[
\boxed{
W_t\mid(W_a=x_a,W_b=x_b)
\sim
\mathcal N\left(
\frac{(b-t)x_a+(t-a)x_b}{b-a},
\varepsilon\frac{(t-a)(b-t)}{b-a}I
\right).
}
\]

В частности, при $a = 0$, $b = 1$ получаем
\[
p(x_t \mid x_0, x_1) = \mathcal N\big((1-t)\,x_0 + t\,x_1,\ \varepsilon\,t(1-t)\,I\big).
\]

\newpage

\section{Второе задание. Цель в лоссе BM}

Пусть $\mathbb{P}$ --- мера на путях процесса $dX_t = \sqrt\varepsilon\, dW_t$, $X_0 = x_0$, а $p(s, y \mid t, x)$ --- его переходная плотность. Хотим определить меру $\mathbb{Q}$, соответствующую условию $X_T = x_1$: $\forall A \in \mathcal{F}_t, \, t < T: \, \mathbb{Q}(A) = \mathbb{P}(A \mid X_T = x_1)$. Применив формулу Байеса, нетрудно получить выражение для правой части:

\begin{equation}\label{eq:cond}
  \mathbb{P}(A \mid X_T = x_1) = \frac{\mathbb{E}_\mathbb{P} \left[ \mathbf{1}_A \ p(T, x_1 \mid t, X_t) \right]}{p(T, x_1 \mid 0, x_0)}
\end{equation}

Введём функцию $h(t, x) = p(T, x_1 \mid t, x) $. Подставив её в \eqref{eq:cond}, получим следующее:

\begin{equation}
  \mathbb{Q}(A) = \mathbb{E}_\mathbb{P} \left[ \mathbf{1}_A \frac{h(t, X_t)}{h(0, x_0)} \right]
\end{equation}

\begin{definition}[Производная Радона --- Никодима]
  \(
\mathbb Q\ll\mathbb P
\implies
\exists Z=\frac{d\mathbb Q}{d\mathbb P}\geq 0:
\quad
\mathbb Q(A)=\int_A Z\,d\mathbb P
\).  $Z$ --- производная Радона --- Никодима.

где \(A\in\mathcal F\), а \(\mathbb Q\ll\mathbb P\) означает:
\(
\mathbb P(A)=0\implies\mathbb Q(A)=0.
\)
\end{definition}

Нетрудно проверить, что если \(
\mathbb Q(A)=\mathbb E_{\mathbb P}[\mathbf 1_A Z],
\)
то $\dfrac{d\mathbb{Q}}{d\mathbb{P}} = Z$. Таким образом, получаем 

\begin{equation}
  Z_t = \frac{d\mathbb{Q}}{d\mathbb{P}} \bigg\vert_{\mathcal{F}_t} = \frac{h(t, X_t)}{h(0, x_0)}
\end{equation}

Переходная плотность $\mathbb{P}$ гауссова, поэтому $h$ известна явно:
\[
h(t, x) = \mathcal N\big(x_1;\, x,\, \varepsilon (T - t) I\big),
\qquad
\nabla_x \log h(t, x) = \frac{x_1 - x}{\varepsilon (T - t)}.
\]
Применим формулу Ито к $h(t, X_t)$ для SDE $dX_t = \sqrt\varepsilon\, dW_t$:

\[
\begin{aligned}
  dh(t, X_t) &= \left(\partial_t h + \frac{\varepsilon}{2} \Delta h\right) dt + \sqrt\varepsilon\, \nabla h^\top dW_t \\&= \sqrt\varepsilon\, \nabla h^\top dW_t,
\end{aligned}
\]
поскольку $h$ как функция начальной точки удовлетворяет обратному уравнению теплопроводности $\partial_t h + \frac{\varepsilon}{2} \Delta h = 0$.

Так как $Z_t = h(t, X_t) / h(0, x_0)$, то
\[
dZ_t = Z_t \sqrt\varepsilon\, \nabla \log{h(t, X_t)}^\top dW_t = Z_t \frac{(x_1 - X_t)^\top}{\sqrt\varepsilon\,(T-t)}\, dW_t.
\]
По теореме Гирсанова процесс $\widetilde W_t = W_t - \int_0^t \frac{x_1 - X_s}{\sqrt\varepsilon\,(T-s)}\, ds$ является винеровским относительно $\mathbb{Q}$, поэтому под $\mathbb{Q}$
\[
dX_t = \frac{x_1 - X_t}{T - t}\, dt + \sqrt\varepsilon\, d\widetilde W_t.
\]
Тот же дрейф можно получить и напрямую из переходной плотности.

\begin{lemma}[Доказательства не будет]
  \[
q(y,s\mid x,t)
=
p(y,s\mid x,t)\frac{h(s,y)}{h(t,x)}.
\]
\end{lemma}

Используя лемму 2.1 для броуновского движения, получим
\begin{equation}\label{eq:step}
\begin{gathered}
  X_{t + \Delta} \mid X_t = x \;\sim\; \mathcal{N}\left( x + \frac{\Delta}{T-t} (x_1 - x), \, \varepsilon \Delta \frac{T - t - \Delta}{T - t} I \right) \quad \text{под } \mathbb{Q}, \\
  \mathbb{E}_\mathbb{Q} \left[ X_{t + \Delta} - X_t \mid X_t = x \right] = \frac{\Delta}{T - t} (x_1 - x).
\end{gathered}
\end{equation}

Поделив \eqref{eq:step} на $\Delta$ и перейдя к пределу $\Delta \downarrow 0$, получим дрейф броуновского моста $(x_1 - x)/(T - t)$. Ковариация приращения, делённая на $\Delta$, стремится к $\varepsilon I$, т.е. коэффициент диффузии остаётся равным $\sqrt\varepsilon$. При $T = 1$ это и есть цель в лоссе bridge matching: $(x_1 - x_t)/(1 - t)$.~$\square$




\end{document}
